
% Evaporation Specific Commands
\newcommand{\BPE}{BPE} % Boiling Point Elevation
\newcommand{\Psol}{P_{\text{sol}}} % Vapor pressure of solution
\newcommand{\Ppure}{P_{\text{pure}}} % Vapor pressure of pure solvent
\newcommand{\xwater}{x_{\text{H}_2\text{O}}} % Mole fraction of water
\newcommand{\gammaAct}{\gamma} % Activity coefficient
\newcommand{\mf}{m_F} % Mass flow rate of feed
\newcommand{\mv}{m_V} % Mass flow rate of vapor
\newcommand{\ml}{m_L} % Mass flow rate of concentrated liquid
\newcommand{\xs}{x_S} % Mass fraction of solute
\newcommand{\Hs}{H_S} % Enthalpy of heating steam
\newcommand{\Hc}{H_C} % Enthalpy of condensate
\newcommand{\HF}{H_F} % Enthalpy of feed
\newcommand{\HV}{H_V} % Enthalpy of vapor
\newcommand{\HL}{H_L} % Enthalpy of concentrated liquid
\newcommand{\Usystem}{U} % Overall heat transfer coefficient for evaporator
\newcommand{\Aheat}{A} % Heat transfer area for evaporator
\newcommand{\DeltaTavail}{\Delta T_{\text{avail}}} % Available temperature difference
\newcommand{\Economy}{E} % Economy of evaporator
\newcommand{\Capacity}{C} % Capacity of evaporator

\section{Introduction to Evaporation Systems}

\subsection{What is an Evaporation System?}
An evaporation system is a specialized type of heat exchanger that facilitates the separation of a solvent (usually water) from a non-volatile solute by vaporizing the solvent. The primary purpose is to concentrate a solution, leading to a concentrated liquid product and a vapor stream. This process involves phase change (boiling) and is crucial in many industrial applications.

\subsection{Applications}
Evaporation systems are widely used in various industries for tasks such as:
\begin{itemize}
    \item Concentrating food products (e.g., fruit juices, milk, sugar solutions).
    \item Chemical industries (e.g., concentrating caustic soda, recovering solvents).
    \item Pharmaceuticals (e.g., concentrating extracts).
    \item Wastewater treatment.
\end{itemize}

\section{Theory and Concepts in Evaporation}

\subsection{Boiling Point Elevation (BPE)}
\begin{itemize}
    \item \textbf{Definition:} The boiling point of a solution containing non-volatile solutes is typically higher than that of the pure solvent at the same pressure. This phenomenon is called boiling point elevation (BPE).
    \item \textbf{Significance:} BPE is crucial in evaporator design because it reduces the effective temperature driving force ($\Delta T$) available for heat transfer. The vapor produced from the solution is at a higher temperature than pure solvent vapor at the same pressure.
    \item \textbf{Dühring Chart:} A graphical tool used to estimate the boiling temperature of an aqueous solution at a given pressure, knowing the boiling temperature of pure water at that same pressure.
        \begin{itemize}
            \item \textbf{Dühring's Rule:} States that the boiling point of a solution varies linearly with the boiling temperature of the pure solvent (e.g., water) at constant solute concentration. This linearity is observed on Dühring charts.
        \end{itemize}
\end{itemize}

\subsection{Vapor-Liquid Equilibrium for Solutions}
\begin{itemize}
    \item \textbf{Raoult's Law:} For ideal solutions, the partial pressure of a solvent (e.g., water) above a solution is proportional to the mole fraction of the solvent in the solution and the vapor pressure of the pure solvent at the same temperature.
    \item \textbf{Activity Coefficient:} For non-ideal solutions (common in industrial evaporation), the actual partial pressure deviates from Raoult's law. This deviation is accounted for by an activity coefficient, which can be determined from experimental data. The vapor pressure of the solution is the partial pressure of the solvent.
\end{itemize}

\subsection{Evaporator Model (Single Effect)}
A single effect evaporator consists of a heat exchanger where steam provides the heating, and a separator where the concentrated liquid product and vapor are separated.
\begin{itemize}
    \item \textbf{Assumptions for Analysis:}
        \begin{enumerate}
            \item The vapor produced is a single component (e.g., pure water vapor).
            \item Vigorous agitation of the heating surface promotes boiling.
            \item A sufficiently large temperature difference ($\Delta T$) is maintained to achieve high evaporation rates.
            \item Boiling point elevation is accounted for.
            \item The overall heat transfer coefficient ($\Uoverall$) is assumed constant.
            \item Negligible heat losses to the surroundings.
            \item No superheat of the vapor (vapor leaves at saturation temperature corresponding to solution pressure).
            \item No crystals formed in the evaporator (assuming concentration below solubility limit).
        \end{enumerate}
    \item \textbf{Balances:} Mass and energy balances are applied to the evaporator system to determine unknown flow rates, temperatures, and heat transfer rates.
\end{itemize}

\subsection{Types of Evaporation Systems}
Evaporators come in various designs, optimized for different applications and solution properties:
\begin{itemize}
    \item \textbf{Single Effect Evaporator:} The simplest type, where the vapor produced is typically condensed and not reused for heating.
    \item \textbf{Multiple Effect Evaporator System:}
        \begin{itemize}
            \item \textbf{Principle:} Reuses the latent heat of the vapor from one effect as the heating steam for the subsequent effect (at a lower pressure). This significantly improves energy efficiency (economy).
            \item \textbf{Configurations:}
                \begin{itemize}
                    \item \textit{Forward-Feed:} Feed enters the first effect and flows sequentially to subsequent effects. Most common for heat-sensitive materials.
                    \item \textit{Backward-Feed:} Feed enters the last effect and flows backward to the first effect. Good for highly viscous solutions.
                    \item \textit{Mixed-Feed/Parallel-Feed:} Combinations to optimize performance for specific needs.
                \end{itemize}
        \end{itemize}
    \item \textbf{Forced-Circulation Evaporator:} Uses a pump to circulate the liquid through the heat exchanger tubes at high velocities, preventing fouling and enhancing heat transfer, especially for viscous liquids or those prone to scaling.
    \item \textbf{Falling-Film Evaporator:} The liquid flows as a thin film down the inside surfaces of vertical tubes. This design is suitable for heat-sensitive materials (short residence time) and for concentrating highly viscous solutions.
    \item \textbf{Short-Tube Vertical Evaporator (Standard/Basket):} Tubes are short and vertical. Boiling occurs inside the tubes, and circulation is driven by density differences.
    \item \textbf{Long-Tube Vertical Evaporator (Rising/Falling Film):} Similar to short-tube but with longer tubes, offering higher heat transfer coefficients and capacity. Can be used for rising film (natural circulation) or falling film (pumped circulation).
\end{itemize}

\section{Performance Measures}

\subsection{Capacity}
The capacity of an evaporator is defined as the number of kilograms of solvent (e.g., water) evaporated per hour.

\subsection{Economy}
The economy of an evaporator is defined as the number of kilograms of solvent evaporated per kilogram of heating steam supplied. Multiple effect evaporators generally have higher economies than single effect evaporators.

 % Start formulas on a new page for clarity
\section{Formulas Summary for Evaporation Systems}

\subsection{Boiling Point Elevation (BPE)}
\subsubsection*{Raoult's Law (for partial pressure of water in solution)}
\begin{equation}
    P_{\text{H}_2\text{O}} = \gammaAct_{\text{H}_2\text{O}} x_{\text{H}_2\text{O}} P_{\text{H}_2\text{O}}^{\text{sat}}
\end{equation}
where: \\
$P_{\text{H}_2\text{O}}$ = Partial pressure of water vapor above the solution (Pa or mmHg) \\
$\gammaAct_{\text{H}_2\text{O}}$ = Activity coefficient of water (dimensionless) \\
$x_{\text{H}_2\text{O}}$ = Mole fraction of water in the liquid phase (dimensionless) \\
$P_{\text{H}_2\text{O}}^{\text{sat}}$ = Saturation vapor pressure of pure water at the solution's temperature (Pa or mmHg)

\subsubsection*{Dühring Chart Interpretation}
The Dühring chart is used to graphically determine the boiling point of a solution at a given pressure (or corresponding boiling temperature of pure water) by aligning known points and interpolating. No explicit formula is typically used for direct calculation from the chart, but it represents the relationship:
$\T_{\text{solution}} = f(\T_{\text{water}}, \text{concentration of solute})$

\subsection{Evaporator Model (Single Effect)}
\subsubsection*{Total Mass Balance}
\begin{equation}
    \dot{m}_F = \dot{m}_V + \dot{m}_L
\end{equation}
where: \\
$\dot{m}_F$ = Mass flow rate of feed solution (kg/s) \\
$\dot{m}_V$ = Mass flow rate of vapor produced (kg/s) \\
$\dot{m}_L$ = Mass flow rate of concentrated liquid product (kg/s)

\subsubsection*{Mass Balance on the Solute}
\begin{equation}
    \dot{m}_F x_{S,F} = \dot{m}_L x_{S,L}
\end{equation}
where: \\
$x_{S,F}$ = Mass fraction of solute in the feed (dimensionless) \\
$x_{S,L}$ = Mass fraction of solute in the concentrated liquid (dimensionless)

\subsubsection*{Overall Energy Balance (assuming negligible heat losses)}
\begin{equation}
    \dot{m}_S H_S + \dot{m}_F H_F = \dot{m}_S H_C + \dot{m}_V H_V + \dot{m}_L H_L
\end{equation}
where: \\
$\dot{m}_S$ = Mass flow rate of heating steam (kg/s) \\
$H_S$ = Enthalpy of heating steam (J/kg) \\
$H_F$ = Enthalpy of feed solution (J/kg) \\
$H_C$ = Enthalpy of condensate from heating steam (J/kg) \\
$H_V$ = Enthalpy of vapor produced (J/kg) \\
$H_L$ = Enthalpy of concentrated liquid product (J/kg) \\
\quad \textbf{Note:} For heating steam, the heat supplied is often $\dot{m}_S (H_S - H_C)$, which is $\dot{m}_S \lambda_S$ if steam condenses as saturated liquid, where $\lambda_S$ is the latent heat of heating steam.

\subsubsection*{Overall Heat Transfer Rate (for evaporator)}
\begin{equation}
    \Qdot = \Uoverall \A_{\text{heat}} \Delta T_{\text{overall}}
\end{equation}
Also, heat transferred from heating steam = heat transferred to solution to vaporize solvent:
\begin{equation}
    \Qdot = \dot{m}_S (H_S - H_C) = \dot{m}_V \lambda_V \quad \text{(simplified, neglecting sensible heat changes in solution)}
\end{equation}
where: \\
$\Qdot$ = Total heat transfer rate in the evaporator (W) \\
$\Uoverall$ = Overall heat transfer coefficient of the evaporator (W/m$^2$$\cdot$K) \\
$\A_{\text{heat}}$ = Heat transfer area of the evaporator (m$^2$) \\
$\Delta T_{\text{overall}}$ = Overall effective temperature difference for heat transfer (K or $^\circ$C). This is typically the temperature of the heating steam minus the boiling temperature of the solution.

\subsection{Multiple Effect Evaporator System}
\subsubsection*{Overall Heat Transfer Rate for each Effect $j$}
\begin{equation}
    \Qdot_j = \Uoverall_j \A_j \Delta T_j
\end{equation}
where $\Delta T_j$ is the temperature difference across effect $j$.

\subsubsection*{Available Temperature Difference}
The total available temperature difference across a multiple effect system is typically the temperature difference between the heating steam inlet to the first effect and the boiling temperature in the last effect (or condenser temperature).
\begin{equation}
    \Delta T_{\text{avail}} = T_{S,1} - T_{V,N} - \sum_{j=1}^{N} \text{BPE}_j
\end{equation}
where: \\
$T_{S,1}$ = Temperature of heating steam in the first effect (K or $^\circ$C) \\
$T_{V,N}$ = Boiling temperature of vapor from the last effect (or condenser temperature) (K or $^\circ$C) \\
$\text{BPE}_j$ = Boiling point elevation in effect $j$ (K or $^\circ$C)

\subsubsection*{Economy (E)}
\begin{equation}
    E = \frac{\text{Total vapor evaporated}}{\text{Total heating steam supplied}}
\end{equation}
(dimensionless)

\subsubsection*{Capacity (C)}
\begin{equation}
    C = \text{Total vapor evaporated per unit time} = \sum_{j=1}^{N} \dot{m}_{V,j}
\end{equation}
(kg/s or kg/hr)
